\section{Lista 1 (05/04/2026)}

\begin{exercise}
	(1.1.4)
	Seja $\Omega = C([0,1])$ com a $\sigma$-álgebra $\mathcal{F} =
	\mathcal{B}(C[0,1])$ (com a topologia induzida
	pela norma do sup). Mostre que
	\[
		\{f \in C([0,1]) : \forall t\; f(t) \geq 0 \} \in \mathcal{F}
	\]
\end{exercise}

\begin{proof}
	% Dado $\alpha \in \R$, seja $g_\alpha \in C([0,1])$ a função
	% constante em $\alpha$.
	Dados $n,m \geq 1$ inteiros, considere as bolas abertas
	\[
		A_{n,m} = B(g_{n/2}, n/2 + 1/m) = \{f \in C([0,1]) : ||f -
		n/2||_\infty < n/2 + 1/m\}
	\]
	onde $g_{n/2}(t) = n/2$ para todo $t$. Notamos, já que toda $f \in C([0,1])$ é limitada, que
	\[
		B_m = \bigcup_{n \geq 1} A_{n,m} = \{f \in C([0,1]) : f > -1/m\}
		\in \mathcal{F}
	\]
	e portanto $\{f \in C([0,1]) : f \geq 0\} = \bigcap_{m\geq 1} B_m
	\in \mathcal{F}$.
\end{proof}

\begin{exercise}
	(1.2.5)
	Seja $n \geq 1$ um número inteiro e considere $\Omega= \{0, 1\}^n$,
	o hipercubo
	de dimensão $n$.
	Para cada $i \in \{1, . . . , n\}$, definimos o evento $A_i =
	\{\omega \in \Omega : \omega(i) = 1\}$. Dadas
	duas probabilidades $\Pr$ e $\Pr '$ em $(\Omega, P(\Omega))$,
	mostre que se $\Pr(B) = \Pr'(B)$ para todos
	conjuntos $B$ dados por interseções de $A_i$, então $\Pr = \Pr'$.
\end{exercise}
\begin{proof}
	% ia repetir a prova de dynkin
	% Sendo $\Omega$ um conjunto finito e já que estamos tomando a
	% $\sigma$-álgebra das partes, segue que
	% para todo $E \in P(\Omega)$
	% \[
	%     \Pr(E) = \sum_{x \in E} \Pr\{x\}
	% \]
	% e o mesmo vale para $\Pr'$. Então basta mostrar que para todo $x
	% \in \Omega$, $\Pr\{x\} = \Pr'\{x\}$.

	% Pela hipótese e pelo princípio da inclusão e exclusão, dado $J
	% \subset [n]$, vale que
	% \begin{align*}
	%     \Pr(\bigcup_{j \in J} A_j) =& \sum_{k=1} (-1)^{k-1}
	%     \sum_{1 \leq i_1 \leq \dots \leq i_k \leq n} \Pr(A_{i_1} \cap
	% \dots \cap A_{i_k})\\
	%     =& \sum_{k=1} (-1)^{k-1}  \sum_{1 \leq i_1 \leq \dots \leq i_k
	% \leq n} \Pr'(A_{i_1} \cap \dots \cap A_{i_k}) = \Pr'(\bigcup_{j \in J} A_j)
	% \end{align*}
	% logo as medidas são idênticas em uniões de $A_i$. Da mesma forma, segue que

	A hipótese nos diz que se $\Pi = \pi(\{A_i\})$ é o $\pi$-sistema
	formado pelos $A_i$'s, então para todo $C \in \Pi$, $\Pr(C) =
	\Pr'(C)$. Segue, pelo teorema de Dynkin, que para qualquer
	$C \in \sigma(\{A_i\})$, temos $\Pr(C) = \Pr'(C)$.

	Vamos mostrar que $P(\Omega) \subset \sigma(\{A_i\})$ e concluir a
	prova. Dado qualquer elemento $w \in \Omega$,
	seja $I(w) = \{i \in [n] : w(i) = 1\}$, temos
	\[
		\{w\} = \bigcap_{i \in I(w)} A_i \cap \bigcap_{j \not \in I(w)}
		A_j^c \quad \in\quad  \sigma(\{A_i\}).
	\]
	Como $\Omega$ é finito, segue que $P(\Omega) \subset \sigma(\{A_i\})$.
\end{proof}

\begin{exercise}
	(1.3.1)
	Dê um exemplo de um $\lambda$-sistema que não seja uma $\sigma$-álgebra.
\end{exercise}
\begin{proof}
	Eu não fiz esse exercício de fato, vi na internet e Bruno sugeriu
	uma construção simples (o João Ferreira
	sugeriu uma construção não trivial muito interessante também!). Seja
	$\Omega = \{0,1,2,3\}$ e considere
	\[
		S = \{\varnothing, \{1,2\}, \{3,4\}, \{1,3\}, \{2,4\}, \{1,2,3,4\}\}.
	\]
	Vamos mostrar que $S$ é um $\lambda$-sistema, mas não é uma
	$\sigma$-álgebra. É claro que $S$ contem $\varnothing$ e é
	fechado por complemento. Agora, as únicas uniões disjuntas possíveis
	em $S$ (descartando o vazio) são
	\[
		\{1,2\} \cup \{3,4\} = \Omega \quad \text{e} \quad \{1,3\} \cup
		\{2,4\} = \Omega
	\]
	portanto, $S$ é fechado por união disjunta e é $\lambda$-sistema.
	Mas $S$ não é uma $\sigma$-álgebra, pois não possui $\{1,2\} \cap
	\{1,3\} = \{1\}$.
\end{proof}

\begin{exercise}
	(1.4.4)
	Dada $X : (\Omega, \mathcal{F}) \to (\mathcal{E}, \Sigma)$ varíavel
	aleatória (ou função mensurável). A $\sigma$-álgebra
	gerada por $x$ é definida por $\sigma(X) := \{X^{-1}(A) : A \in
	\Sigma\}$. Seja $X:\R^2 \to \R$ dada por $X(w_1,w_2) = \max(w_1,w_2)$.
	Descreva $\sigma(X)$.
\end{exercise}

\begin{proof}
	Não sei exatamente o quão bem quer que eu descreva. Dado $A \in
	\mathcal{B}(\R)$, temos
	\[
		X^{-1}(A) = \{(w_1,w_2) \in \R^2 : \max(w_1,w_2) \in A\} =
		(\{w_1 \leq w_2\} \cap (\R \times A)) \cup (\{w_2 \leq w_1\} \cap
		(A \times \R))
	\]
	logo $\sigma(X)$ é coleção desses conjuntos onde $A \in
	\mathcal{B}(\R)$. Ela não é $\mathcal{B}(\R^2)$,
	pois para todo conjunto não vazio $X^{-1}(A) \in \sigma(X)$, temos
	que $A$ não é vazio e portanto a
	semi-reta $\{(a, t): t \leq a\}$ pertence a $X^{-1}(A)$. Como é
	claramente falso que todo conjunto de $\mathcal{B}(\R^2)$
	possui uma semireta infinita, segue que $\sigma(X) \neq \mathcal{B}(\R^2)$.
\end{proof}

\begin{exercise}
	(1.4.6)
	Seja $X : [0,1] \to \{0,1\}$ dada por $X = \id[A]$ para $A$
	mensurável em $[0,1]$. Mostre que $X_*\Pr = \text{Ber}(p)$ para
	algum $p \in [0,1]$. Calcule o valor de $p$.
\end{exercise}
\begin{proof}
	É claro que
	\[
		X_*\Pr(S) = \Pr(\id[A]^{-1}(S)) =
		\begin{cases}
			1 & \text{se}\; S = \{0,1\}\\
			\Pr(A) & \text{se}\; S = \{1\}\\
			1 - \Pr(A) & \text{se}\; S = \{0\}\\
			0 & \text{se}\; S = \varnothing\\
		\end{cases}
	\]
	portanto, $X_*\Pr = \text{Ber}(\Pr(A))$.
\end{proof}

\begin{exercise}
	(1.4.7)
	Sejam $X,Y$ variáveis aleatórias tais que $X$ é nula quase
	certamente. Mostre que $X+Y$ tem a mesma distribuição de $Y$.
\end{exercise}
\begin{proof}
	Basta notar que, dado $E$ mensurável na reta,
	\begin{align*}
		\Pr(X + Y \in E) &= \Pr((X + Y \in E) \land (X = 0 \lor X \neq 0))\\
		&= \Pr((Y \in E) \land X = 0) + \Pr((X + Y \in E) \land X \neq 0) =
		\Pr(Y \in E).
	\end{align*}
\end{proof}

\begin{exercise}
	(2.2.3 e 2.2.4)
	Seja $X : ([0,1], \mathcal{B}(\R)) \to (S^1,\mathcal{B}(\C))$ dada
	por $X(t) = \exp(2i\pi t)$. Definimos
	\[
		\mathcal{U}_{S^1} := X_*\Pr
	\]
	onde $\Pr = \text{Unif}([0,1])$. Mostre que $\mathcal{U}_{S^1}$ é
	invariante por
	rotação e não é absolutamente contínua em relação a medida
	de Lebesgue em $\C$.
\end{exercise}

\begin{proof}
	Vamos provar a invariância por rotaçao. Dado $E \subset S^1$
	mensurável e $e^{2\pi i\theta} \in S^1$, queremos
	que $X_*\Pr(E) = X_*\Pr(e^{2\pi i \theta}E)$. Sejam
	\[
		A = \{t \in [0,1]: e^{2\pi i t} \in E\} \quad \text{ e } \quad B =
		\{t \in [0,1]: e^{2\pi i t} \in e^{2\pi i \theta}E\},
	\]
	basta ter que $m(A) = m(B)$. Usando $\{a\}_\text{frac}$ para denotar
	a parte fracionária do real $a$, note que (trocando $t - \theta$ por $u$),
	\[
		B = \{t \in [0,1] : e^{2\pi i (t - \theta)} \in E\} = \{ \{u +
		\theta\}_\text{frac} : e^{2\pi i u} \in E\}
	\]
	mas então $m(B) = m(\{A + \theta\}_\text{frac})$ e pela invariância
	por translação da medida de Lebesgue $m(\cdot)$,
	segue que $m(B) = m(A)$. Mais concretamente,
	\begin{align*}
		m(\{A + \theta\}_\text{frac}) &= \int_0^1 \id[x \in \{A +
		\theta\}_\text{frac}] dx = \int_0^1 \id[\{x - \theta\}_\text{frac} \in A] dx\\
		&= \int_0^{\theta} \id[x - \theta + 1 \in A]dx +
		\int_{\theta}^{1}\id[x - \theta \in A] dx\\
		&= \int_{1-\theta}^{1} \id[z \in A] dz + \int_{0}^{1-\theta} \id[u \in A] du
	\end{align*}
	onde a segunda igualdade vale, pois dados $x,y,\theta \in [0,1)$, $x
	= \{y + \theta\}_\text{frac}$ se e somente se $\{x - \theta\}_\text{frac} = y$.
	E a última igualdade vale quando fizemos a substituição $z = x -
	\theta + 1$ e $u = x - \theta$ nas últimas duas integrais.

	% \begin{align*}
	% 	\mathcal{U}_{S^1}(e^{2\pi i\theta} E) &= \Pr( \{t \in [0,1] : e^{2i\pi
	% 	t} \in e^{2\pi i\theta} E\})\\
	% 	&= \Pr(\{t \in [0,1] : e^{2i\pi (t - \theta)} \in E\})\\
	% 	&= \Pr(\{ \{u + \theta\}_\text{frac} \in [0,1] : e^{2i\pi u} \in E\})
	% 	% = \mathcal{U}_{S^{-1}}(E)
	% \end{align*}
	% onde a terceira igualdade segue da invariância da medidade de
	% Lebesgue por translação. De maneira mais concreta,
	% dado $U \subset [0,1]$ e $\alpha \in [0,1)$, então $m(U) = m(\{ \{u +
	% \alpha\}_\text{frac} : u \in U\})$, pois
	% \begin{align*}
	% 	m(\{U + a\}_\text{frac}) &= \int_{0}^1 \id[\{x - a\} \in U] dx \\
	% 	&=\int_{0}^{1-a} \id[x + a \in U] dx + \int_{1-a}^1 \id [x + a -
	% 1 \in U] dx\\
	% 	&= \int_{a}^{1} \id[y \in U] dy + \int_{0}^{a} \id[z \in U] dz = m(U)
	% \end{align*}
	% onde fizemos a substituição $y = x + a$ e $z = x + a - 1$.

	A medida não é absolutamente contínua em relação à Lebesgue no plano
	pois $m(S^1) = 0$ enquanto $\mathcal{U}_{S^1}(S^1) = 1$.
\end{proof}

\begin{exercise}
	(2.3.2)
	Calcule $F_{\text{Exp}(\lambda)}$ ($F_X$ é a notação da função
	cumulativa usada pelo Augusto).
\end{exercise}
\begin{proof}
	\[
		F_{\text{Exp}(\lambda)}(t) = \int_{0}^{t} \lambda e^{-\lambda x} dx
		= 1 - e^{-\lambda t}
	\]
\end{proof}

\begin{exercise}
	(2.3.3)
	Mostre, usando o Teorema da Extensão de Caratheodory, que se $F:\R \to [0,1]$ é
	uma função que satisfaz:
	\begin{enumerate}[label=(\roman*)]
		\item $\lim_{x\to -\infty} F(x) = 0$, $\lim_{x\to \infty} F(x) = 1$
		\item $F$ é monótona não decrescente
		\item $F$ é contínua à direita e possui limite à esquerda (càdlag).
	\end{enumerate}

	Então existe uma única medida de probabilidade $\mu$ em $(\R,
	\mathcal{B}(\R))$ tal que $F$ é função cumulativa de $\mu$.
\end{exercise}

Para esse exercício vou enunciar o Teorema de Extensão de Caratheodory e
seguir a prova do Durret.
\begin{theorem}
	\label{trm:caratheodory_extension}
	(Extensão de Caratheodory). Seja $\mathcal{G} \subset
	\mathcal{P}(\Omega)$ uma álgebra
	de conjuntos em $\Omega$ e suponha que $\mu : \mathcal{G} \to \R^+$
	satisfaça que:
	para toda família finita ou enumerável $(A_i)_{i \in I}$ de
	conjuntos disjuntos de $\mathcal{G}$
	tal que $\bigsqcup_{i \in I} A_i \in \mathcal{G}$,
	temos que $\mu(\bigsqcup_{i \in I} A_i) = \sum_{i \in I} \mu(A_i)$.
	Então existe uma medida $\hat{\mu} : \sigma(\mathcal{G}) \to \R^+$
	tal que $\hat{\mu}(A) = \mu(A)$ para todo $A \in \mathcal{G}$.

	% Seja $\mathcal{S}$ uma semi-álgebra e $\mu$ definida em
	% $\mathcal{S}$ com $\mu(\varnothing) = 0$. Suponha que $\mu$ satisfaz:
	% \begin{enumerate}[label=(\roman*)]
	% 	\item Se $S, (S_i)_{i=1}^{N} \in \mathcal{S}$ com $S = \bigsqcup_{i
	% 		= 1}^N S_i$ então $\mu(S) = \sum_{i=1}^N \mu(S_i)$.
	% 	\item Se $S, (S_i)_{i=1}^{\infty} \in \mathcal{S}$ com $S =
	% 		\bigsqcup_{i = 1}^\infty S_i$ então $\mu(S) \leq
	% \sum_{i=1}^\infty \mu(S_i)$.
	% \end{enumerate}
	% Então $\mu$ tem uma única extensão $\hat{\mu}$ que é uma medida na
	% álgebra $A(\mathcal{S})$ gerada por $S$. Se $\hat{\mu}$
	% é $\sigma$-finita então existe uma única extensão $\nu$ que é uma
	% medida em $\sigma(\mathcal{S})$.
\end{theorem}

\begin{proof}
	(Do exercício)
	Seja $\mathcal{G}$ a álgebra gerada pelo conjunto dos intervalos
	semiabertos $(a,b]$ com
	$-\infty \leq a \leq b \leq \infty$. Lembramos (de um curso usual de
	medida) que todo elemento $S \in \mathcal{G}$
	se escreve como $\bigsqcup_{i = 1}^n (a_i,b_i]$ para alguma
	sequência $(a_i,b_i]$ de intervalos semiabertos disjuntos.

	Definimos $\mu: \mathcal{G} \to \R^+$ dada por 
	\[
		\mu(S) = \begin{cases}
			F(b) - F(a) & \text{ se } S = (a,b]\\
			\sum_{i = 1}^{N} F(b_i) - F(a_i) & \text{ se } S = \bigsqcup_{i=1}^{N} (a_i,b_i]
		\end{cases}
	\]
	Vamos verificar que essa definição de fato faz sentido quando escrevemos nossos
	conjuntos de maneiras distintas.
	% [\ref{trm:caratheodory_extension}].

	(Escrevendo um intervalo como união disjunta) Note que se $(a,b] =
	\bigsqcup_{i=1}^N
	(a_i,b_i]$ então devemos ter $a_1 = a$, $b_N = b$ e
	$b_i = a_{i+1}$ para todo $1 \leq i \leq N-1$ . Portanto,
	\[
		\mu((a,b]) = F(b) - F(a) = \sum_{i = 1}^{N} F(b_i) - F(a_i) = \mu(\bigsqcup_{i=1}^{N} (a_i,b_i])
	\]
	e portanto a definição faz sentido (fazemos o abuso notacional
	$F(-\infty) = 0$ e $F(\infty) = 1$).

	(Escrevendo uniões de forma distinta) Seja $A = \bigsqcup_{i=1}^{N}
	(a_i,b_i]$. Se $A$ é da forma $(a,b]$ já provamos
	o resultado. Então vamos supor (tomando intervalos maximais)
	que $A$ é da forma $\bigsqcup_{j = 1}^M (x_j, y_j]$ onde cada $(x_j,y_j]$ é uma componente conexa de $A$.
	Segue que para cada $i \in [N]$, existe um único $j \in [M]$ tal que $(a_i,b_i] \subseteq (x_j,y_j]$ e
	além do mais, sendo $I(j) := \{i \in [N] : (a_i,b_i] \subset (x_j,y_j]\}$, vale que
	\[
	(x_j,y_j] = \bigsqcup_{i\in I(j)} (a_i,b_i].
	\]
	Temos então, pelo resultado anterior,
	\[
		\mu(\bigsqcup_{j = 1}^M (x_j, y_j]) = \sum_{j = 1}^{M} F(y_j) - F(x_j) = 
		\sum_{j=1}^{M} \sum_{i \in I(j)} F(b_i) - F(a_i) = \sum_{i = 1}^N \mu((a_i,b_i]) = \mu(\bigsqcup_{i=1}^N (a_i,b_i])
	\]
	como queríamos mostrar. Ou seja a definição faz sentido.

	Vamos ver que $\mu$ satisfaz as condições de Caratheodory [\ref{trm:caratheodory_extension}]. Como toda união 
	finita de elemento de $\mathcal{G}$ é uma união finita de intervalos semiabertos, já verificamos a igualdade para 
	uniões disjuntas finitas. Basta então considerar uniões infinitas enumeráveis $S = \bigsqcup_{i \geq 1} A_i \in \mathcal{G}$,
	que, reindexando e abrindo os $A_i$'s, podem ser escrito da forma:
	\[
	S = \bigsqcup_{i = 1}^\infty A_i = \bigsqcup_{i = 1}^{\infty} \bigsqcup_{j = 1}^{n_i} (a_{i,j},b_{i,j}]
	= \bigsqcup_{k = 1}^{\infty} (a_k, b_k].
	\]
	Precisamos que $\mu(S) = \sum_{k = 1}^{\infty} F(b_k) - F(a_k)$. Primeiramente, fica claro que 
	$\mu(S) \geq \sum_{k = 1}^{\infty} F(b_k) - F(a_k)$ pois, para todo $N$, temos 
	\begin{align*}
		\mu(S) &= \mu\bigg(\bigsqcup_{k = 1}^{N}(a_k, b_k] \sqcup \bigg(S \setminus \bigsqcup_{i = 1}^{N} (a_k, b_k] \bigg) \bigg)\\
		&= \sum_{k = 1}^N F(b_k) - F(a_k) + \mu\bigg(S \setminus \bigsqcup_{i = 1}^{N} (a_k, b_k]\bigg)\\
		&\geq \sum_{k = 1}^N F(b_k) - F(a_k).
	\end{align*}
	
	Falta a outra desigualdade $\mu(S) \leq \sum_{k = 1}^{\infty} F(b_k) - F(a_k)$. Recorreremos
	da hipótese de continuidade a direita. Inicialmente, seja $S = (a,b]$ com $-\infty < a \leq b < \infty$
	e suponha novamente $S = \bigsqcup_{i=1}^{\infty} (a_i,b_i]$. Dado $\eps > 0$
	arbitrário, vamos mostrar que
	\[
		\mu((a,b]) = F(b) - F(a) \leq \eps + \sum_{i = 1}^{\infty}
		\mu((a_i,b_i]) = \eps/2 + \sum_{i = 1}^{\infty} \mu((a_i,b_i]) +
		\frac{\eps}{2^{i+1}}
	\]
	Por continuidade a direita, podemos tomar $\delta > 0$ tão pequento
	tal que $F(a+\delta) - F(a) < \eps/2$.
	Da mesma forma, para cada $i$, tomamos $\delta_i > 0$ com $F(b_i +
	\delta_i) < F(b_i) + \eps/2^{i+1}$. Segue que
	\[
		(a+ \delta, b] \subset [a + \delta, b] \subset \bigcup_{i=1}^\infty
		(a_i, b_i + \delta_i),
	\]
	como o segundo conjunto é compacto, possui subcobertura $1 \leq i_1
	< \dots <i_M$ finita e vale
	\[
		(a + \delta, b] \subset \bigcup_{j=1}^{M} (a_{i_j}, b_{i_j} +
		\delta_{i_j}) \subset \bigcup_{j=1}^{M} (a_{i_j}, b_{i_j} + \delta_{i_j}].
	\]
	Temos
	\begin{equation}
		\label{exer:eq:cebola}
		\mu((a+\delta, b]) = F(b) - F(a + \delta) \leq \sum_{j=1}^{M}
		F(b_{i_j} + \delta_{i_j}) - F(a_{i_j}) \leq \sum_{j=1}^{\infty}
		F(b_j) - F(a_j) + \frac{\eps}{2^{i+1}}
	\end{equation}
	substituindo,
	\[
		\mu((a,b]) = \mu((a + \delta, b]) + F(a + \delta) - F(a) \leq \eps
		+ \sum_{i = 1}^{\infty} \mu((a_i,b_i])
	\]
	como queríamos mostrar.

	Usamos nessa demonstração (em [\ref{exer:eq:cebola}]) uma
	propriedade usual de pre-medidas em álgebras que se $S \subset
	\bigcup_{i=1}^N S_i$ então
	vale que $\mu(S) \leq \sum_{i = 1}^{N} \mu(S_i)$. (É o argumento da cebola).

	Ainda temos que mostrar o resultado para $a,b$ potencialmente
	$-\infty$ e $\infty$. Mas isso segue da existência do limite da $F$.
	Se $(A,B] \subset (a,b] \subset \bigsqcup_{i=1}^{\infty} (a_i, b_i]$
	com $-\infty < A \leq B < \infty$, vimos que
	$\mu((A,B]) \leq \sum_{i = 1}^{\infty} \mu((a_i,b_i])$. Fazendo $A$
	ou $B$ tender para $-\infty$ ou $+\infty$ nos dá o resultado geral.

	Se $S$ é da forma $S = \bigsqcup_{j = 1}^{M} (x_j,y_j] = \bigsqcup_{i = 1}^{\infty} (a_i, b_i]$,
	segue, do mesmo argumento de antes, que chamando $I(j) := \{i \in \N : (a_i,b_i]  \subset (x_j, y_j] \}$
	(esses agora podem ser infinitos), achamos 
	\[
		\mu(S) = \mu(\bigsqcup_{j = 1}^{M} (x_j, y_j]) = \sum_{j = 1}^M \sum_{i \in I(j)} \mu((a_i,b_i])
		= \sum_{i = 1}^{\infty} F(b_i) - F(a_i).
	\]

	Por Caratheodory [\ref{trm:caratheodory_extension}], temos uma
	medida na $\sigma$-álgebra gerada pelo semi-abertos, que é $\mathcal{B}(\R)$. Como a álgebra 
	dos semiabertos é um $\pi$-sistema, essa medida é única pelo teorema de Dynkin.
\end{proof}

\begin{exercise}
	(2.3.12)
	Se $X \in L^1(P)$ e $\Pr(X \geq x) = \Pr(X \leq -x)$ para todo $x$.
	Então $\Ex X = 0$.
\end{exercise}
\begin{proof}
	Basta escrever (notando que $\Ex[X \id[X = 0]] = 0$) que
	\begin{align*}
		\Ex X &= \Ex[X \id[X \geq 0]] + \Ex[X\id[X \leq 0]] - \E[X \id[X = 0]]\\
		&= \Ex[X \id[X \geq 0]] - \Ex[-X \id[X \leq 0]]\\
		&= \int_{0}^{\infty} \Pr(X\id[X \geq 0] \geq t)dt -
		\int_{0}^{\infty} \Pr(-X\id[X \leq 0] \geq t)dt\\
		&= \int_{0}^{\infty} \Pr(X \geq t)dt - \int_{0}^{\infty} \Pr(X \leq -t)dt\\
		&= \int_{0}^{\infty} \Pr(X \geq t) - \Pr(X \leq -t) dt = 0
	\end{align*}
\end{proof}
